\section*{結言}
この日のXでは、AI開発をどこまで協調して安全側へ寄せるかという政策・ガバナンス論が、製品発表と並ぶ主要テーマになった。安全性協議を進めるとの報道と、競争や各社責任を重視して協調的減速に距離を置くとの報道が同時に流れ、フロンティア企業間でも単一の方向へ収束していない状況が映し出された。DeepMind Instituteの立ち上げも、AGIの社会的影響を能力開発と並行して研究する動きとしてこの流れに重なる。

一方、MistralとMozillaの提携、Grok BuildやGrok Bot、Paper2Agentなどでは、モデル単体よりもブラウザ、エージェント、MCP、実行環境との接続が引き続き前面に出ている。資金調達やセキュリティ事案には未確認情報も残るため、翌日以降は公式発表の有無と、ライブ企画・研究成果が実際の利用や再現へどうつながるかが継続観測点となる。
